\section{Application of Program Outcomes}
Table~\ref{tab:po} maps each Program Outcome and Program Specific Outcome addressed by this case study to its concrete implementation. Every mapping demonstrates how foundational engineering knowledge, modern tools, and professional ethics are exercised within the project.

{\small
\begin{longtable}{@{}>{\centering\arraybackslash}p{1.8cm} p{13.2cm}@{}}
    \caption{Mapping of Program Outcomes and Program Specific Outcomes to the Case Study.}\label{tab:po}\\
    \toprule
    \textbf{Outcome} & \textbf{Outcome Description and Concrete Relevance to the Case Study} \\
    \midrule
    \endfirsthead
    \toprule
    \textbf{Outcome} & \textbf{Outcome Description and Concrete Relevance to the Case Study} \\
    \midrule
    \endhead
\textbf{PO1} & \textbf{Engineering Knowledge:} Applies principles of mathematics, computer science, and engineering fundamentals to formalize, model, and analyze the system architecture, mathematical constraints, and foundational algorithms. In this case study (Predictive Data Analytics for Student Academic Performance Evaluation), mathematical formulations of statistical risk functions, time-series decompositions, and information metrics are applied to Predictive Data Analytics for Student Academic Performance Evaluation. \\
    \addlinespace
\textbf{PO2} & \textbf{Problem Analysis:} Systematically identifies, formulates, and investigates complex computational challenges, diagnosing architectural bottlenecks, concurrency anomalies, latency constraints, and integrity failures through rigorous analytical techniques. In this case study (Predictive Data Analytics for Student Academic Performance Evaluation), systematic problem formulation resolves missing data anomalies, class imbalance, and high-dimensional noise in analytical datasets. \\
    \addlinespace
\textbf{PO3} & \textbf{Design/Development of Solutions:} Designs scalable system architectures, robust database schemas, object hierarchies, or automated pipelines meeting specified operational needs with proper consideration for efficiency, maintainability, and safety. In this case study (Predictive Data Analytics for Student Academic Performance Evaluation), a multi-stage data architecture spanning ingestion, transformation, predictive modeling, and executive dashboarding is designed. \\
    \addlinespace
\textbf{PO4} & \textbf{Conduct Investigations of Complex Problems:} Uses research-based knowledge and experimental methodologies, including synthesis of empirical benchmarks, stress testing, and profiling, to draw valid technical conclusions. In this case study (Predictive Data Analytics for Student Academic Performance Evaluation), rigorous empirical validation using 10-fold cross-validation, statistical significance testing, and SHAP explainability provides experimental rigor. \\
    \addlinespace
\textbf{PO5} & \textbf{Modern Tool Usage:} Selects and applies state-of-the-art software frameworks, modeling tools, database management engines, profiling utilities, and development environments to deliver high-performance solutions. In this case study (Predictive Data Analytics for Student Academic Performance Evaluation), modern data science environments including Pandas, Scikit-Learn, LightGBM, Plotly, and Tableau are proficiently utilized. \\
    \addlinespace
\textbf{PO6} & \textbf{The Engineer and Society:} Evaluates contextual consequences of technical deployments on institutional workflows, user accessibility, social well-being, and legal compliance. In this case study (Predictive Data Analytics for Student Academic Performance Evaluation), societal impacts on student success, healthcare equity, and commercial transparency are evaluated and prioritized. \\
    \addlinespace
\textbf{PO12} & \textbf{Life-Long Learning:} Recognizes the need for, and develops the preparation and ability to engage in independent lifelong learning in the broadest context of technological evolution. In this case study (Predictive Data Analytics for Student Academic Performance Evaluation), analytical competencies developed through this case study prepare students for the rapid evolution of enterprise business intelligence. \\
    \addlinespace
\textbf{PSO1} & \textbf{Program-Specific (AI/ML Core):} Leverages advanced intelligent computational techniques, algorithmic optimization, predictive analytics, and automated decision-making workflows to enhance system capabilities. In this case study (Predictive Data Analytics for Student Academic Performance Evaluation), advanced predictive modeling, classification scorecards, and econometric forecasting directly reinforce AI/ML expertise. \\
    \addlinespace
\textbf{PSO2} & \textbf{Program-Specific (Applied Engineering):} Designs and implements production-grade end-to-end software pipelines integrating intelligent modules with robust backend infrastructure and enterprise services. In this case study (Predictive Data Analytics for Student Academic Performance Evaluation), end-to-end software integration bridging raw data lakes with production dashboard APIs fulfills applied engineering mandates. \\
    \addlinespace
    \bottomrule
\end{longtable}}